\documentclass[11pt,letterpaper]{article}
\usepackage[margin=0.85in,headheight=16pt]{geometry}
\usepackage[T1]{fontenc}
\usepackage[utf8]{inputenc}
\usepackage{lmodern,microtype}
\usepackage{amsmath,amssymb,mathtools,booktabs,longtable,tabularx,array}
\usepackage[dvipsnames]{xcolor}
\usepackage{hyperref,fancyhdr,titlesec,enumitem,graphicx,tcolorbox}
\usepackage{lastpage}
\definecolor{ink}{HTML}{14263B}
\definecolor{sea}{HTML}{167C8A}
\definecolor{gold}{HTML}{C49747}
\definecolor{pale}{HTML}{EEF5F6}
\definecolor{graytext}{HTML}{516275}
\hypersetup{colorlinks=true,linkcolor=ink,urlcolor=sea,bookmarksopen=true,pdfauthor={Independent Educational Engineering Suite}}
\titleformat{\section}{\Large\bfseries\color{ink}}{\thesection}{0.7em}{}[\vspace{-0.3em}\color{gold}\titlerule]
\titleformat{\subsection}{\large\bfseries\color{sea}}{\thesubsection}{0.6em}{}
\setlist[itemize]{leftmargin=1.3em,itemsep=2pt,topsep=3pt}
\setlist[enumerate]{leftmargin=1.5em,itemsep=2pt,topsep=3pt}
\pagestyle{fancy}
\fancyhf{}
\fancyhead[L]{\small\sffamily\color{graytext}SECURITY ANALYSIS / ENGINEERING SUITE}
\fancyhead[R]{\small\sffamily\color{graytext}\textbf{\thepage}/\pageref{LastPage}}
\renewcommand{\headrulewidth}{0.3pt}
\renewcommand{\footrulewidth}{0pt}
\newtcolorbox{methodbox}[1][]{colback=pale,colframe=sea,boxrule=0.55pt,arc=1mm,title={\bfseries Method note},fonttitle=\sffamily,#1}
\newtcolorbox{sourcebox}[1][]{colback=white,colframe=gold,boxrule=0.55pt,arc=1mm,title={\bfseries Source boundary},fonttitle=\sffamily,#1}
\newcommand{\term}[1]{\textbf{\color{ink}#1}}
\newcommand{\src}{\textsf{\color{sea}[SOURCE]}}
\newcommand{\ext}{\textsf{\color{gold}[EXTENSION]}}
\newcommand{\gap}{\textsf{\color{red}[SOURCE GAP]}}
\newcommand{\answer}[1]{\par\smallskip\noindent\textbf{Check:} #1\par}
\newcommand{\cover}[3]{%
\begin{titlepage}\pagecolor{ink}\color{white}\vspace*{1.15in}
{\sffamily\large\color{gold}INDEPENDENT TECHNICAL DOCUMENTATION SUITE\par}\vspace{0.4in}
{\Huge\sffamily\bfseries #1\par}\vspace{0.15in}
{\Large\sffamily #2\par}\vspace{0.65in}
\color{gold}\rule{\textwidth}{2pt}\par\vspace{0.4in}
{\large #3\par}\vfill
{\small\sffamily Based on a critical reading of Graham and Dodd's \emph{Security Analysis}, Seventh Edition (2023). Original educational analysis; not affiliated with the publisher.\par}
\end{titlepage}\nopagecolor\color{black}}
\newcommand{\front}[3]{\cover{#1}{#2}{#3}\tableofcontents\clearpage
\begin{sourcebox}\src\ means a concept is attributable to the relevant chapter or seventh-edition contributor, summarized independently rather than reproduced. \ext\ means a modern method or engineering recommendation prepared for this suite, not an original Graham--Dodd rule. Historical examples and standards should not be treated as current market requirements. All numerical examples are synthetic, educational, and not securities recommendations.\end{sourcebox}}
\setlength{\parskip}{5pt}\setlength{\parindent}{0pt}
\begin{document}
\front{Investment Valuation}{Quantitative Methods and Model Controls}{Worked DCF, multiples, NCAV, liquidation, bonds, duration, credit and conversion}

\section{Model governance and conventions}
\ext\ This quantitative handbook is an original engineering extension; it does not attribute DCF, WACC, duration, Monte Carlo simulation or portfolio expected shortfall to the historical chapters. Values are hypothetical. Model inputs must specify fiscal period, reporting currency, units, accounting policy, source, transformation, and reviewer.
\begin{methodbox}\textbf{Dimensional control:} keep price per share ($\$/\text{share}$), market capitalisation ($\$$), enterprise value ($\$$), and debt market value ($\$$) distinct. For rates, use decimal values inside formulas (e.g., 9\%=0.09).\end{methodbox}

\section{Discounted cash flow: FCFF and enterprise value}
\subsection{Forecast model}
\[
FCFF_t=EBIT_t(1-T_t)+D\!\&A_t-CapEx_t-\Delta NWC_t,\quad
EV_0=\sum_{t=1}^{n}\frac{FCFF_t}{(1+r)^t}+\frac{TV_n}{(1+r)^n}.
\]
A perpetuity-growth terminal value $TV_n=FCFF_{n+1}/(r-g)$ requires $r>g$ and economically sustainable long-run reinvestment. Terminal exit multiples and perpetuity assumptions are alternative models and should not be combined or chosen solely to reach a target price.

\subsection{Reproducible example}
Assume FCFF in years 1--3 of \$80m, \$90m, \$100m; discount rate 10\%; terminal growth 3\%; net debt \$400m; diluted shares 100m.
\[
TV_3=\frac{100(1.03)}{0.10-0.03}=1{,}471.43\text{m},
\]
\[
EV=\frac{80}{1.10}+\frac{90}{1.10^2}+\frac{100+1{,}471.43}{1.10^3}
\approx1{,}327.74\text{m}.
\]
Thus equity value $\approx\$927.74$m and value/share $\approx\$9.28$. Always check terminal value share of EV, terminal reinvestment realism, and sensitivity to rates.

\subsection{WACC}
For debt $D$, equity $E$, tax rate $T$, equity required return $r_e$ and pre-tax debt cost $r_d$,
\[
WACC=\frac{E}{D+E}r_e+\frac{D}{D+E}r_d(1-T).
\]
This textbook-style expression is appropriate only under matching definitions and stable financing assumptions. Distressed firms, tax-loss positions, preferred capital, changing leverage and country risk demand an adapted method. \emph{Never combine nominal cash flows with real discount rates or mismatch currencies.}

\section{Reverse DCF and scenario boundaries}
Instead of assuming value, solve for revenue growth, terminal margins, reinvestment or duration of excess returns implied by market price. A simple reverse growth model uses a target $EV_{\text{market}}$ and finds $g$ such that
\[
0=\sum_t \frac{FCFF_t(g)}{(1+r)^t}+\frac{FCFF_{n+1}(g)}{(r-g_T)(1+r)^n}-EV_{\text{market}}.
\]
Keep terminal growth $g_T$ independent of near-term operating growth and constrain both economically. \textbf{Review:} compare implied economics with market size, capacity, competitive response and return on incremental invested capital.

\section{Owner earnings, normalized EPS and multiples}
\[
NOPAT=EBIT(1-T),\quad ROIC=\frac{NOPAT}{\text{average invested capital}},\quad
\text{P/E}=\frac{P}{EPS}.
\]
Valuation multiples require normalized financials and compatible numerators/denominators: EV/EBITDA, EV/EBIT, and equity P/E are \emph{not} interchangeable. Current ROIC above WACC implies value-creating incremental growth only if new capital can earn an adequate incremental return.

\subsection{Illustration: quality-adjusted multiple}
Firm A has EV \$900m and maintainable EBIT \$90m; EV/EBIT = 10.0x. Firm B has EV \$840m and maintainable EBIT \$70m; EV/EBIT = 12.0x. B is more highly valued despite lower EV. Investigate ROIC, reinvestment needs, business durability and debt before concluding which is expensive.

\section{Asset, net-current-asset and liquidation methods}
\[
BVPS=\frac{\text{book common equity}}{N},\qquad
NCAVPS=\frac{CA-TL}{N},
\]
where $CA$ is current assets and $TL$ is \emph{all} liabilities, not merely current liabilities. Example: cash 30, receivables 40, inventory 50, total liabilities 70, shares 10m (all amounts \$m). $NCAV=30+40+50-70=50$m and $NCAVPS=\$5$.

For liquidation, suppose cash 30 realizes 100\%, receivables 40 realize 85\%, inventory 50 realizes 50\%, and net disposal costs equal 7. Then proceeds after liabilities equal $30+34+25-70-7=\$12$m, or \$1.20/share. This differs markedly from NCAV and illustrates the fragility of apparent asset coverage.

\section{Sum of the parts and capitalization waterfall}
For each business segment $i$, calculate segment enterprise value $EV_i$ on consistent economics. Then
\[
E=\sum_i EV_i+\text{surplus cash}-\text{gross debt}-\text{priority claims}-\text{holdco costs}.
\]
Avoid counting intersegment profits twice or adding investment stakes already reflected in consolidated EBITDA. Distribute realizable equity to share classes according to actual instruments, not simply share count.

\section{Fixed-income present value, yield and sensitivity}
For annual coupon $C$, face $F$, annual yield $y$, and maturity $n$,
\[
P=\sum_{t=1}^{n}\frac{C}{(1+y)^t}+\frac{F}{(1+y)^n}.
\]
Illustration: 3-year annual 5\% coupon on \$1,000 face discounted at 7\%:
\[
P=\frac{50}{1.07}+\frac{50}{1.07^2}+\frac{1050}{1.07^3}\approx\$947.51.
\]
Yield-to-maturity assumes scheduled payments and reinvestment conventions implicit in the yield interpretation; it is not an expected realized return under credit risk.

\subsection{Duration and DV01}
\[
D_{Mac}=\frac{1}{P}\sum_t t\frac{CF_t}{(1+y)^t},\qquad
D_{mod}=\frac{D_{Mac}}{1+y},
\quad DV01\approx P D_{mod}(10^{-4}).
\]
Here DV01 applies to a 1 basis point increase in annual yield under the simplified parallel-yield assumption. For large shifts, include convexity and spread/rate curve distinctions. Callability, floating coupons, and embedded options require option-adjusted measures.

\section{Credit expected loss and distressed recovery}
A simple one-period model is $EL=PD\times LGD\times EAD$ where probability of default, loss-given-default, and exposure-at-default correspond to \emph{the same horizon}. Illustration: $PD=4\%$, $LGD=60\%$, $EAD=\$10$m yields $EL=\$240{,}000$. Expected loss does not describe the magnitude of a single adverse realization or account for correlated defaults.

For recovery, split operating going-concern EV from liquidation realization and allocate through an issuer-specific waterfall. Record discount to face price versus expected cash recovery, timing, court costs and legal priority. Scenario weighting must not conceal an unacceptable worst-case funding need.

\section{Convertible claims and warrants}
For $F=\$1{,}000$ convertible face amount and conversion price $K=\$40$, $q=25$ shares. If stock price $S=\$48$, conversion parity is $\$1{,}200$. At convertible price $P_c=\$1{,}260$, parity premium is $1260/1200-1=5\%$. The stock can fall while bond-floor protection also deteriorates because issuer credit risk may worsen. \ext\ For simple warrant expiry payoff, $\max(0,S_T-K)$ per warrant ignores dilution and financing; explicitly correct those in issuer-level analyses.

\section{Sensitivity and model review matrix}
\begin{center}\small\begin{tabularx}{\textwidth}{p{0.28\textwidth}X X}\toprule Variable & Base perturbation & Failure mode inspected \\
\midrule WACC/required return & $\pm 200$ bps & Long-duration terminal valuation \\
Revenue or margins & $\pm 15\%$ cash flow & Operating leverage and customer risk \\
Terminal growth & $0\%$ to $3\%$ & Implied perpetual reinvestment \\
Inventory realization & $30\%$ to $80\%$ & Liquidation shortfall \\
Refinancing spread & $+300$ bps & Interest cover and maturity wall \\
Equity dilution & $+5\%$ shares & Per-share claim degradation \\
FX depreciation & $-25\%$ local currency & Investor-currency value preservation \\
\bottomrule\end{tabularx}\end{center}

\section{Verification checklist and reproducible output}
\begin{enumerate}
\item Source and date for each input; no unsupported hard-coded market quote.
\item Discounted versus undiscounted units and timing consistent; horizon at year-end or midyear documented.
\item Enterprise-to-equity bridge ties to balance sheet; debt and minority interests not subtracted twice.
\item Stock counts are diluted consistently with contingent claims and option exercise proceeds.
\item Downside values reconcile to potential loss of principal and forced-sale liquidity.
\item Model checks: $r>g$, positive denominator, no circular references without controlled solution, and independent reviewer sign-off.
\end{enumerate}

\end{document}
